% Preliminaries figure: a CMOS cell = PUN (PMOS) + PDN (NMOS), NAND2 example
\begin{tikzpicture}[font=\small, >=Stealth,
    lbl/.style={font=\scriptsize}]
\ctikzset{transistors/scale=0.62}

%% ===================== CMOS schematic (left) =====================
% supply rails
\draw[thick] (-0.30,4.15) -- (2.55,4.15); \node[lbl,left] at (-0.30,4.15) {$V_\mathrm{DD}$};
\draw[thick] (-0.30,-0.35) -- (2.55,-0.35); \node[lbl,left] at (-0.30,-0.35) {GND};

% --- PUN: two PMOS in parallel (pull-up) ---
\node[pmos, yscale=-1] (p1) at (0.35,3.35) {};
\node[pmos, yscale=-1] (p2) at (1.95,3.35) {};
\draw (p1.S) -- (0.35,4.15);  \draw (p2.S) -- (1.95,4.15);
\draw (p1.D) -- (0.35,2.55) -- (1.95,2.55) -- (p2.D);
\node[lbl,left=1pt] at (p1.G) {$a$};
\node[lbl,left=1pt] at (p2.G) {$b$};

% --- output node ZN ---
\draw (1.15,2.55) -- (1.15,2.20); \filldraw (1.15,2.55) circle (1.1pt);
\node[lbl,right] at (2.05,2.55) {$\mathrm{ZN}=\lnot(a\,b)$};

% --- PDN: two NMOS in series (pull-down) ---
\node[nmos] (n1) at (1.15,1.55) {};
\node[nmos] (n2) at (1.15,0.45) {};
\draw (n1.D) -- (1.15,2.20);
\draw (n1.S) -- (n2.D);
\draw (n2.S) -- (1.15,-0.35);
\node[lbl,left=1pt] at (n1.G) {$a$};
\node[lbl,left=1pt] at (n2.G) {$b$};

% role labels
\node[lbl, align=center, blue!55!black] at (3.55,3.35)
    {\textbf{PUN}\\PMOS\\pull \emph{up}\\when $f{=}1$};
\node[lbl, align=center, red!60!black] at (3.55,1.00)
    {\textbf{PDN}\\NMOS\\pull \emph{down}\\when $f{=}0$};

%% ===================== truth table (right) =====================
\node[anchor=north west] at (4.75,4.35) {%
  \small\setlength{\tabcolsep}{4pt}\renewcommand{\arraystretch}{1.05}
  \begin{tabular}{cc|c|l}
    $a$ & $b$ & $\mathrm{ZN}$ & \\ \hline
    0 & 0 & 1 & \scriptsize on \\
    0 & 1 & 1 & \scriptsize on \\
    1 & 0 & 1 & \scriptsize on \\
    1 & 1 & 0 & \scriptsize off \\
  \end{tabular}};
\node[lbl, anchor=north west, align=left, text width=3.0cm] at (4.70,1.35)
    {on-patterns $\Pon$: PUN conducts.\\[2pt]
     off-pattern $\Poff$: PDN conducts.};

\end{tikzpicture}
